\subsection{Objective}

Task 6 requires segmenting the characters so that the different characters
are separated as clearly as possible and each one is individually labelled.
The input is the binary image selected in Task~4, in which the ten characters
of \texttt{HD44780A00} already appear as ten separate bright regions on a dark
background.  The segmentation therefore has three steps: remove the residual
noise, label every remaining connected component, and order the ten labelled
characters from left to right.

\subsection{Segmentation Method}

\subsubsection{Connected-component labelling}

Foreground components were labelled with the classic two-pass
union-find algorithm~\cite{gonzalez2018}, implemented in \texttt{common.py}.  In the first pass the
image is scanned row by row; every foreground pixel takes the smallest label
among its already-visited 8-neighbours, and the corresponding labels are
merged in a union-find structure.  In the second pass every label is replaced
by the root of its equivalence class.  The result is a label image in which
each 8-connected foreground region has a unique label.

\subsubsection{Noise removal by area}

Task~4 reported 126 components, of which 10 are characters and 116 are texture
residues.  The largest noise component has 224 pixels while the smallest
character has $2{,}719$ pixels, so an area criterion removes all noise without
touching any character:
\begin{equation}
  \text{keep component } C \iff \lvert C\rvert \geq 300 .
  \label{eq:task6-area}
\end{equation}
The threshold is deliberately placed in the wide gap between the two
populations, which makes the choice insensitive to small image changes.

\subsubsection{Left-to-right ordering}

The surviving components are sorted by their leftmost column so that label $1$
is the leftmost character and label $10$ the rightmost, matching the reading
order of \texttt{HD44780A00}.

\subsection{Results}

\cref{fig:task6-labeled} shows the ten labelled characters over the grayscale
sub-image, and \cref{fig:task6-boxes} adds the bounding box of every label.
The boxes are $79$--$88$ pixels wide and $114$--$131$ pixels high, so the
characters are remarkably uniform in size.  \cref{fig:task6-crops} shows the
ten extracted character crops that are saved for the classification stage.

\begin{figure}[!htbp]
  \centering
  \includegraphics[width=\textwidth]{task6_labeled.png}
  \caption{Segmented characters labelled $1$--$10$ in reading order.}
  \label{fig:task6-labeled}
\end{figure}

\begin{figure}[!htbp]
  \centering
  \includegraphics[width=\textwidth]{task6_boxes.png}
  \caption{Bounding boxes of the ten segmented characters.}
  \label{fig:task6-boxes}
\end{figure}

\begin{figure}[!htbp]
  \centering
  \includegraphics[width=\textwidth]{task6_crops.png}
  \caption{The ten extracted character crops (labels $1$--$10$).}
  \label{fig:task6-crops}
\end{figure}

\begin{table}[!htbp]
  \centering
  \caption{Segmentation result.  Bounding boxes are given in sub-image pixel
  coordinates.}
  \label{tab:task6-components}
  \begin{tabular}{clccc}
    \toprule
    Label & Glyph & Bounding box $(x_{0},y_{0})$--$(x_{1},y_{1})$ &
    Size ($w\times h$) & Area (px) \\
    \midrule
    1 & \texttt{H} & (8, 19)--(91, 138) & $84\times120$ & 3836 \\
    2 & \texttt{D} & (101, 10)--(185, 140) & $85\times131$ & 4149 \\
    3 & \texttt{4} & (190, 18)--(269, 144) & $80\times127$ & 3485 \\
    4 & \texttt{4} & (280, 19)--(361, 134) & $82\times116$ & 3269 \\
    5 & \texttt{7} & (372, 10)--(450, 135) & $79\times126$ & 2719 \\
    6 & \texttt{8} & (453, 19)--(540, 135) & $88\times117$ & 4336 \\
    7 & \texttt{0} & (544, 16)--(626, 138) & $83\times123$ & 3859 \\
    8 & \texttt{A} & (630, 19)--(713, 132) & $84\times114$ & 3103 \\
    9 & \texttt{0} & (716, 15)--(801, 135) & $86\times121$ & 3980 \\
    10 & \texttt{0} & (808, 13)--(889, 136) & $82\times124$ & 3997 \\
    \bottomrule
  \end{tabular}
\end{table}

\subsection{Why Component Labelling Rather than Projection Splitting}

The horizontal gaps between adjacent bounding boxes are only
$9$, $4$, $10$, $10$, $2$, $3$, $3$, $2$, and $6$ pixels, so several
characters almost touch.  A purely projection-based splitter would therefore
be fragile.  The effect is visible when the same labelling is applied to the
global-Otsu binary from Task~4: there the \texttt{D} and first \texttt{4}
merge, the \texttt{8} and \texttt{0} merge, and the two final \texttt{0}s
merge, leaving only 7 groups instead of 10
(\cref{fig:task6-comparison}).  The adaptive binary keeps the strokes
separated, so the simple component-labelling segmentation already returns
exactly one region per character, and no further de-merging or split heuristics
are needed.

\begin{figure}[!htbp]
  \centering
  \includegraphics[width=\textwidth]{task6_segmentation_comparison.png}
  \caption{Character groups found on the global-Otsu binary (7 groups) and on
  the selected adaptive binary (10 groups).  Merged neighbours are visible as
  boxes that span two characters.}
  \label{fig:task6-comparison}
\end{figure}

\subsection{Analysis}

The segmentation is exact for this row: after noise removal exactly ten
components remain, and reading the components left to right reproduces
\texttt{H D 4 4 7 8 0 A 0 0}.  The result inherits two properties from the
binarization stage.  First, the \texttt{8} remains one component even though
its bottom loop is partly open (Task~5), because the stroke itself is still
connected; area filtering does not break it further.  Second, the three
\texttt{0}s have slightly different areas ($3{,}859$, $3{,}980$, $3{,}997$
pixels) because of local stroke-thickness variation, which shows why a fixed
area interval rather than a single expected size is used to accept a
component.

The cleaned binary is saved as \path{task_6/outputs/task6_cleaned_binary.png},
and one padded grayscale and binary crop per character is saved under
\path{task_6/outputs/chars/} together with
\path{task_6/outputs/task6_components.json}.  These crops are the natural
input of the CNN and non-CNN classifiers in Task~7, whose ground-truth label
order is exactly the known string \texttt{HD44780A00}.

\FloatBarrier
